% preamble.tex - POPL port of the Polymorphic Type Slicing macros.
% Loaded AFTER \documentclass{acmart}. acmart already provides amsmath, amssymb,
% amsthm, graphicx, xcolor, hyperref, booktabs, geometry, natbib, and predefined
% theorem/lemma/proposition/corollary/conjecture/definition/example environments.
% Dropped vs the dissertation: geometry, fancyhdr, hyperref, biblatex, the lualatex
% font setup (directlua emdash + \setmainfont), caption, marginnote, all
% mechanisation glyphs/marginpars (now no-ops; see end of file).

% acmart loads amsmath but not amssymb; we need \triangleq, \blacktriangleright,
% \Box, \boxdot, \nRightarrow, ... throughout.
% acmart's newtxmath already defines \Bbbk; clear it so amssymb won't abort.
\let\Bbbk\relax
\usepackage{amssymb}

% Save \sqsubseteq before semantic redefines it (needed for tikz-cd labels).
\let\sqsubseteqOriginal\sqsubseteq
\let\sqsupseteqOriginal\sqsupseteq

\usepackage{semantic}        % \inference

% --- Fix semantic.sty inference-rule spacing under LuaLaTeX + unicode-math ---
% infernce.sty spaces the rule bar with `\vskip 4\fontdimen8\textfont3`:
% fontdimen8 of the legacy cmex10 math-extension font is the default rule
% thickness (0.4pt). acmart >= v2.16 under LuaLaTeX loads unicode-math, making
% \textfont3 the OpenType math font (Libertinus), where that fontdimen is
% ~6.8pt -- inserting ~27pt of glue above AND below every rule bar. Patch the
% \vskip's (and the nested-inference raise) back to the classic 0.4pt value.
\makeatletter
\newdimen\sem@defaultrulethick
\sem@defaultrulethick=0.4pt
\patchcmd{\@inferenceBack}
  {\vskip 4\fontdimen8\textfont3}{\vskip 4\sem@defaultrulethick}
  {}{\PackageWarning{preamble}{semantic \string\inference\space patch 1 failed}}
\patchcmd{\@inferenceBack}
  {\vskip 4\fontdimen8\textfont3}{\vskip 4\sem@defaultrulethick}
  {}{\PackageWarning{preamble}{semantic \string\inference\space patch 2 failed}}
\patchcmd{\@processInference}
  {\advance\dimen3 by \fontdimen8\textfont3}
  {\advance\dimen3 by \sem@defaultrulethick}
  {}{\PackageWarning{preamble}{semantic \string\inference\space patch 3 failed}}
\makeatother
% ----------------------------------------------------------------------------
\usepackage{stmaryrd}        % \llbracket, \rrparenthesis, ...
\providecommand{\Circle}{\bigcirc} % wasysym incompatible with lualatex
\usepackage{mathtools}

% acmart >= v2.16 loads unicode-math with Libertinus Math, whose stretchy
% horizontal braces (U+23DE/U+23DF) render as solid black bars under
% \underbrace/\overbrace. Pull just those two glyphs from Latin Modern Math.
\makeatletter
\@ifpackageloaded{unicode-math}{%
  \setmathfont{latinmodern-math.otf}[range={"023DE,"023DF},Scale=MatchLowercase]%
}{}
\makeatother

\usepackage{soul}             % \hl
\usepackage[most]{tcolorbox}
\usepackage[normalem]{ulem}

\usepackage{tabularx}
\usepackage{arydshln}         % dashed col separators
\usepackage{float}
\usepackage{subcaption}

\usepackage{etoolbox}
\usepackage{needspace}
% NOTE: acmart errors on redefining sectioning commands ("An attempt to redefine
% \subsection detected"), so the \pretocmd{\subsection}{\needspace{...}} hack from
% the dissertation is removed for the POPL version.

\usepackage{cleveref}
% Abbreviated section cross-references: \cref -> "Sec. 3" mid-sentence,
% \Cref -> "Section 3" at sentence start. Sub(sub)sections inherit these.
\crefname{section}{Sec.}{Secs.}
\Crefname{section}{Section}{Sections}

% ---------------------------------------------------------------------------
% Colours
% ---------------------------------------------------------------------------
\definecolor{commentgreen}{RGB}{2,112,10}
\definecolor{weborange}{RGB}{255,165,0}
\definecolor{frenchplum}{RGB}{129,20,83}

% Vibrant overrides for dvipsnames slice/highlight colours
\definecolor{BrickRed}{RGB}{210,35,40}
\definecolor{NavyBlue}{RGB}{30,80,225}
\definecolor{OliveGreen}{RGB}{35,165,45}
\definecolor{Mulberry}{RGB}{215,55,175}

% Universal-property diagram convention: assumed (Class 2), induced (Class 3)
\colorlet{assumed}{BurntOrange!80!black}
\colorlet{induced}{TealBlue!45!black}

% ---------------------------------------------------------------------------
% Review-edit highlighting (every reworded/new span is wrapped for the author).
% \edited{...} for inline spans; editedblock for new paragraphs / display math.
% ---------------------------------------------------------------------------
% lua-ul \highLight survives line breaks and arbitrary macros (\cref, footnotes)
% under lualatex, unlike soul's \hl.
\usepackage{luacolor}
\usepackage{lua-ul}
\newcommand{\edited}[1]{\highLight[cyan!22]{#1}}
\newtcolorbox{editedblock}{breakable, enhanced, colback=cyan!6, colframe=cyan!45,
  boxrule=0.3pt, left=4pt, right=4pt, top=3pt, bottom=3pt, sharp corners}

% ---------------------------------------------------------------------------
% Code / slice text fonts
% ---------------------------------------------------------------------------
\newcommand{\asmref}[1]{\textcolor{assumed}{\textbf{#1}}}
\newcommand{\synsrc}[1]{\textcolor{NavyBlue!60!gray}{#1}}
\newcommand{\anasrc}[1]{\textcolor{OliveGreen!60!gray}{#1}}
% Heavy mono for highlighted type-source spans. The dissertation used
% SourceCodePro-Black via fontspec; here we use bold typewriter to avoid a
% hard font dependency under acmart.
\newcommand{\heavymono}{\bfseries\ttfamily}
\newcommand{\synsrcS}[1]{\textcolor{blue!85!black}{{\heavymono #1}}}
\newcommand{\anasrcS}[1]{\textcolor{OliveGreen!75!black}{{\heavymono #1}}}
\newcommand{\strN}[1]{\textcolor{black!85}{#1}}

% Fold glyph: alias for \gap.
\newcommand{\fold}{\gap}

% ---------------------------------------------------------------------------
% Lattice ops and precision in a muted colour
% ---------------------------------------------------------------------------
\definecolor{latticecolor}{HTML}{8055B5}
\let\OLDsqcap\sqcap
\let\OLDsqcup\sqcup
\let\OLDsqsubset\sqsubset
\let\OLDsqsupset\sqsupset
\def\sqcap{\mathbin{{\color{latticecolor}\OLDsqcap}}}
\def\sqcup{\mathbin{{\color{latticecolor}\OLDsqcup}}}
\def\sqsubseteq{\mathrel{{\color{latticecolor}\sqsubseteqOriginal}}}
\def\sqsupseteq{\mathrel{{\color{latticecolor}\sqsupseteqOriginal}}}
\def\sqsubset{\mathrel{{\color{latticecolor}\OLDsqsubset}}}
\def\sqsupset{\mathrel{{\color{latticecolor}\OLDsqsupset}}}
\newcommand{\heytingto}{\mathbin{{\color{latticecolor}\to}}}
\newcommand{\heytingminus}{\mathbin{{\color{latticecolor}\setminus}}}
\newcommand{\latticetop}{{\color{latticecolor}\top}}
\newcommand{\latticebot}{{\color{latticecolor}\bot}}

% ---------------------------------------------------------------------------
% Highlighting
% ---------------------------------------------------------------------------
\newcommand{\hlc}[2][yellow]{{%
    \colorlet{foo}{#1}%
    \sethlcolor{foo}\hl{#2}}%
}
\newcommand{\error}[1]{\hlc[pink!50]{#1}}
\newcommand{\hlcmaths}[2][yellow!50]{\colorbox{#1}{$\displaystyle #2$}}

\newcommand{\dyn}{{\color{teal}\texttt{?}}}          % dynamic type
\newcommand{\gap}{{\color{teal}\Box}}
\newcommand{\hole}[1][]{{\color{teal}\llparenthesis} #1 {\color{teal}\rrparenthesis}} % hole
\newcommand{\cmark}{{\color{blue}\Circle}}            % context mark
\newcommand{\type}[1]{\llbracket #1 \rrbracket}       % type interpretation
\newcommand{\id}{\mathrm{id}}
\newcommand{\code}[1]{\texttt{#1}}

% Well-formedness: Delta |- tau wf
\newcommand{\wf}[2][\Delta]{#1 \vdash #2\;\mathsf{wf}}

% Synthesis / analysis judgements
\newcommand{\synthesis}[3][\Delta;\,\Gamma]{#1 \vdash #2 {\color{BrickRed}\ \Uparrow}\ #3}
\newcommand{\analysis}[3][\Delta;\,\Gamma]{#1 \vdash #2 {\color{BlueViolet}\ \Downarrow}\ #3}

% Rule-name labels
\newcommand{\synr}[1]{\textnormal{{\color{BrickRed}$\Uparrow$}\texttt{#1}}}
\newcommand{\anar}[1]{\textnormal{{\color{BlueViolet}$\Downarrow$}\texttt{#1}}}
\newcommand{\cnsr}[1]{\textnormal{$#1$}}

\newcommand{\typeassignment}[3][\Delta;\Gamma]{#1 \vdash #2 : #3}

% Slice judgements
\newcommand{\synthesisslice}[4][\Gamma]{#1 \vdash #2 {\color{BrickRed}\ \Rightarrow}\ #3 \mid #4}
\newcommand{\analysisslice}[5][\Gamma]{#1; #2 \vdash #3 {\color{BlueViolet}\ \Leftarrow}\ #4 \mid #5}

% Term-minimal slice judgement (used in the Optimising section and the supplement).
\providecommand{\tmrule}[7][\Delta;\,\Gamma]{%
  #1 \vdash #2 \mathbin{{\color{BrickRed}\Uparrow}} #3
  \mathbin{{\color{BrickRed}\blacktriangleleft}} #4
  \;{\color{BrickRed}\rightsquigarrow}\; #5
  \mathbin{{\color{BrickRed}\Uparrow}} #6
  \,\dashv\, #7%
}

% Casts
\newcommand{\scast}[2]{{\color{blue}\langle} #1 {\color{blue}\Rightarrow} #2 {\color{blue}\rangle}}
\newcommand{\scasterror}[2]{{\color{red}\langle} #1 {\color{red}\Rightarrow} \dyn {\color{red}\nRightarrow} #2 {\color{red}\rangle}}
\newcommand{\scastcast}[3]{{\color{blue}\langle} #1 {\color{blue}\Rightarrow} #2 {\color{blue}\Rightarrow} #3 {\color{blue}\rangle}}

% Type matching
\newcommand{\funmatch}{\blacktriangleright_\rightarrow}
\newcommand{\groundmatch}{\blacktriangleright_{\text{ground}}}
\newcommand{\forallmatch}{\blacktriangleright_\forall}

% Elaboration judgements
\newcommand{\elaborationSynthesis}[5][\Gamma]{#1 \vdash #2 {\color{BrickRed}\ \Uparrow\ } #3 \leadsto #4 \dashv #5}
\newcommand{\elaborationAnalysis}[6][\Gamma]{#1 \vdash #2 {\color{BlueViolet}\ \Downarrow\ } #3 \leadsto #4 : #5 \dashv #6}
\newcommand{\elaborationSynthesisSlice}[6][\Gamma]{#1 \vdash #2 {\color{BrickRed}\ \Rightarrow\ } #3 \mid #4 \leadsto #5 \dashv #6}
\newcommand{\elaborationAnalysisSlice}[9][\Gamma]{#1;\ #2 \vdash #3 {\color{BlueViolet}\ \Leftarrow\ } #4 \mid #5 \leadsto #6 : #7 \mid #8 \dashv #9}

% Context classification
\newcommand{\ctxclass}[5][\Delta;\,\Gamma]{#1 \vdash #2\;\mathbf{at}\;#3\;\triangleright\;#4\;[\,#5\,]}

% Marking
\newcommand{\marksto}{\looparrowright}
\newcommand{\synmark}[4][\Delta;\,\Gamma]{#1 \vdash #2 \marksto #3 {\color{BrickRed}\ \Uparrow}\ #4}
\newcommand{\anamark}[4][\Delta;\,\Gamma]{#1 \vdash #2 \marksto #3 {\color{BlueViolet}\ \Downarrow}\ #4}
\newcommand{\mctxclass}[6][\Delta;\,\Gamma]{#1 \vdash #2 \marksto #3\;\mathbf{at}\;#4\;\triangleright\;#5\;[\,#6\,]}
\newcommand{\chk}[1]{\check{#1}}

% Mark wrappers
\definecolor{markcolor}{RGB}{170,40,60}
\newcommand{\markopen}{{\color{markcolor}\llparenthesis}}
\newcommand{\markclose}{{\color{markcolor}\rrparenthesis}}
\newcommand{\markincon}[3][\cdot]{{\markopen #1 \markclose}_{#2 \not\sim #3}}
\newcommand{\markarrow}[1][\cdot]{{\markopen #1 \markclose}_{\blacktriangleright \Rightarrow}}
\newcommand{\marksum}[1][\cdot]{{\markopen #1 \markclose}_{\blacktriangleright +}}
\newcommand{\markprod}[1][\cdot]{{\markopen #1 \markclose}_{\blacktriangleright \times}}
\newcommand{\markforall}[1][\cdot]{{\markopen #1 \markclose}_{\blacktriangleright \forall}}
\newcommand{\marklam}[1][\cdot]{{\markopen #1 \markclose}_{\sim \Rightarrow}}
\newcommand{\markprodsim}[1][\cdot]{{\markopen #1 \markclose}_{\sim \times}}
\newcommand{\markvar}[1]{{\color{markcolor}\langle} #1 {\color{markcolor}\rangle}\!\Uparrow}

\newcommand{\markbox}[1]{%
  \tikz[baseline=(N.base)]{%
    \node[draw=BrickRed!85, dashed, line width=0.7pt, rounded corners=2pt,
          inner sep=2pt, fill=BrickRed!6, anchor=base] (N) {\strut$#1$};%
  }%
}

\DeclareRobustCommand{\squiggle}[1]{%
  {\color{red}\uwave{\mbox{\normalcolor #1}}}%
}

\newcommand{\slist}[1]{%
  \tikz[baseline=(N.base)]{%
    \node[draw=OliveGreen!25, dashed, line width=0.4pt, rounded corners=1.5pt,
          inner sep=1.5pt, fill=OliveGreen!4, anchor=base] (N) {\strut$#1$};%
  }%
}

\newcommand{\queryblock}[2]{%
  \begin{center}
  \small
  \begin{tabular}{r@{\;}c@{\;}l}
    {\color{NavyBlue}synthesis query}  & $\upsilon\;=$ & $#1$ \\
    {\color{OliveGreen}analysis query} & $\psi\;=$ & $#2$ \\
  \end{tabular}
  \end{center}
}

% Structural variables (from formalism)
\newcommand{\C}{\mathcal{C}}
\newcommand{\Cs}{c}
\newcommand{\p}{d}
\newcommand{\F}{\mathcal{F}}
\newcommand{\f}{f}
\newcommand{\Sb}{\mathcal{S}}

\DeclareMathAlphabet{\mathdcal}{U}{dutchcal}{m}{n}
\newcommand{\A}{\mathdcal{A}}

\newcommand{\ground}[1][\tau]{#1\text{ ground}}
\newcommand{\final}[1][d]{#1\text{ final}}
\newcommand{\val}[1][d]{#1\text{ val}}
\newcommand{\boxedval}[1][d]{#1\text{ boxedval}}
\newcommand{\indet}[1][d]{#1\text{ indet}}

\newcommand{\contextualsub}[1][d/u]{\llbracket #1 \rrbracket}

% Polymorphism
\newcommand{\Forall}[2]{\forall #1.\ #2}
\newcommand{\TLam}[2]{\Lambda #1.\ #2}
\newcommand{\TApp}[2]{#1\ @#2}
\newcommand{\ForallSlice}[2]{\forall #1.\ #2}
\newcommand{\gapforall}{{\color{teal}\forall\Box.\Box}}

% Labelled tuples
\newcommand{\lbl}[2]{\mathtt{#1} : #2}
\newcommand{\tuplety}[1]{\{#1\}}

% ---------------------------------------------------------------------------
% Theorem environments. acmart predefines theorem/lemma/proposition/corollary/
% conjecture/definition/example; we add the ones it lacks.
% ---------------------------------------------------------------------------
\makeatletter
\theoremstyle{acmplain}
\@ifundefined{theorem}{\newtheorem{theorem}{Theorem}[section]}{}
\@ifundefined{lemma}{\newtheorem{lemma}[theorem]{Lemma}}{}
\@ifundefined{proposition}{\newtheorem{proposition}[theorem]{Proposition}}{}
\@ifundefined{corollary}{\newtheorem{corollary}[theorem]{Corollary}}{}
\@ifundefined{conjecture}{\newtheorem{conjecture}[theorem]{Conjecture}}{}
\@ifundefined{counterexample}{\newtheorem{counterexample}[theorem]{Counterexample}}{}
\theoremstyle{acmdefinition}
\@ifundefined{definition}{\newtheorem{definition}[theorem]{Definition}}{}
\@ifundefined{example}{\newtheorem{example}[theorem]{Example}}{}
\theoremstyle{acmplain}
\@ifundefined{remark}{\newtheorem{remark}[theorem]{Remark}}{}
\makeatother

% Master-theorem sub-numbering: produces "Theorem N.M.k"
\newcounter{subthm}
\newcommand{\masterthmnum}{??}
\renewcommand{\thesubthm}{\masterthmnum.\arabic{subthm}}
\newcommand{\mastercases}{%
  \xdef\masterthmnum{\thetheorem}%
  \setcounter{subthm}{0}%
}
\newenvironment{thmcase}[1][]{%
  \refstepcounter{subthm}%
  \par\medskip\noindent{\bfseries Theorem~\thesubthm}%
  \ifx\\#1\\\else\ (\textit{#1})\fi%
  {\bfseries.}\itshape\quad\ignorespaces
}{%
  \par\medskip
}

\newtcolorbox{recipe}[1]{
  colback=blue!3, colframe=blue!40, fonttitle=\bfseries,
  title={Recipe: #1}, sharp corners, boxrule=0.5pt
}
\newtcolorbox{originalerror}[1][]{
  colback=red!3, colframe=red!40, fonttitle=\bfseries,
  title={Error (Original)#1}, sharp corners, boxrule=0.5pt
}
\newtcolorbox{majortheorembox}{
  colback=green!3, colframe=OliveGreen!60, sharp corners, boxrule=0.5pt, enhanced
}
\newtcolorbox{counterexamplebox}{
  colback=orange!3, colframe=orange!50, sharp corners, boxrule=0.5pt, enhanced
}

% Multi-slice highlighting
\newcommand{\sliceA}[1]{\textbf{\underline{\textcolor{BrickRed}{#1}}}}
\newcommand{\sliceB}[1]{\underline{\textcolor{NavyBlue}{#1}}}
\newcommand{\sliceC}[1]{\underline{\textcolor{OliveGreen}{#1}}}
\newcommand{\sliceD}[1]{\underline{\textcolor{Mulberry}{#1}}}

% Highlighted-box slice notation
\newcommand{\sli}[2][NavyBlue]{%
  \tikz[baseline=(N.base)]{%
    \node[draw=#1!60, dashed, line width=0.4pt, rounded corners=1.5pt,
          inner sep=1.5pt, fill=#1!12, anchor=base] (N) {\strut$#2$};%
  }%
}
\newcommand{\slr}[1]{%
  \tikz[baseline=(N.base)]{%
    \node[draw=BrickRed!80, line width=0.8pt, rounded corners=2pt,
          inner sep=3pt, fill=none, anchor=base] (N) {\strut$#1$};%
  }%
}
\newcommand{\fullslice}[2]{%
  \begin{tikzpicture}[baseline=(t.base), every node/.style={inner sep=3pt}]
    \node[draw, line width=0.9pt, rounded corners,
          label={[font=\scriptsize\itshape]above:assumptions}] (a) {$#1$};
    \node[below=6pt of a] (t) {$#2$};
  \end{tikzpicture}%
}
\newcommand{\gapbox}{\boxdot}

% tikz
\usepackage{tikz-cd}
\usetikzlibrary{positioning,calc,fit,backgrounds,decorations.pathmorphing,shapes.geometric,arrows.meta}
\tikzcdset{arrows={->}, diagrams={column sep=large, row sep=large}}

% Lightweight footnote-block environment retained from the dissertation.
\newsavebox{\fnbox}
\newenvironment{fn}
  {\begin{lrbox}{\fnbox}\footnotesize\begin{minipage}{\linewidth}}
  {\end{minipage}\end{lrbox}\footnote{\usebox{\fnbox}}}

% ---------------------------------------------------------------------------
% Mechanisation markers and marginpars: REMOVED for the POPL version (a separate
% mechanisation supplement is provided). Defined as no-ops so any residual call in
% the copied sources still compiles.
% ---------------------------------------------------------------------------
\newcommand{\yo}{}
\newcommand{\no}{}
\newcommand{\po}{}
\newcommand{\iy}{}
\newcommand{\ino}{}
\newcommand{\ipart}{}
\newcommand{\mech}[2]{}
\newcommand{\mechyes}[1]{}
\newcommand{\mechno}[1]{}
\newcommand{\mechpartial}[1]{}
\newcommand{\mechfile}[2][]{}
\newcommand{\mechbox}[1]{}
\newcommand{\sectionlegend}{}
% Raw \marginnote / \marginpar no-ops (signature [opt]{mand}[opt]).
\makeatletter
\renewcommand{\marginpar}[2][]{}
\newcommand{\marginnote}[2][]{\@ifnextchar[\@mn@gobble\relax}
\def\@mn@gobble[#1]{}
\makeatother
